\SourceLeaf{283}%
\begin{center}
\begin{minipage}[t]{5.7cm}\centering $\begin{array}{rcc}
&&a\\
1.8&+&\\
&&b\\
2.9&-&\\
&&a''\\
3.18&-&\\
&&c\\
4.17&+&\\
&&a^{\mathrm{VI}}\\
5.12&+&\\
&&b^{\mathrm{V}}\\
6.21&+&\\
&&d\\
7.14&-&\\
&&b^{\mathrm{VIII}}\\
8.1&-&\\
&&a'\\
9.2&+&\\
&&b'\\
10.19&-&\\
&&a^{\mathrm{IV}}\\
11.16&+&\\
\end{array}$\end{minipage}
\hfill
\begin{minipage}[t]{5.7cm}\centering $\begin{array}{rcc}
&&b^{\mathrm{IV}}\\
12.5&-&\\
&&a^{\mathrm{VII}}\\
13.22&+&\\
&&b^{\mathrm{VII}}\\
14.7&+&\\
&&d'\\
15.20&-&\\
&&b'''\\
16.11&-&\\
&&a^{\mathrm{V}}\\
17.4&-&\\
&&c'\\
18.3&+&\\
&&a'''\\
19.10&+&\\
&&b''\\
20.15&+&\\
&&d''\\
21.6&-&\\
&&b^{\mathrm{VI}}\\
22.13&-&\\
\end{array}$\end{minipage}
\end{center}
\[
(a) = \,.\, -b+c-d
\]
\[
(b) = +a\,.\, +c-d
\]
\[
(c) = -a-b\,.\, -d
\]
\[
(d) = +a+b+c\,.
\]

\Needspace{20\baselineskip}\begin{center}[6.]\end{center}

The proof of the rule given above rests on a different division
of the figure into several closed figures, none of which intersects
itself, although one may intersect another. These are the
figures that form successively as the line is first traced.
Select only the half of the crossings that are encountered
for the second time, in our example
